\documentclass[11pt]{article}
\usepackage{amsmath,amssymb,amsthm}
\usepackage{booktabs}
\usepackage[margin=1in]{geometry}
\newtheorem{theorem}{Theorem}
\newtheorem{lemma}[theorem]{Lemma}
\title{Problem 763: Amoeba Colonies 3D}
\author{Project Euler Solutions}
\date{}
\begin{document}
\maketitle
\section{Problem Statement}
An amoeba at $(0,0,0)$ divides into 3 offspring at $(x+1,y,z)$, $(x,y+1,z)$, $(x,y,z+1)$ if those cells are empty. After $N$ divisions, $2N+1$ amoebas exist. Different division orders may give same arrangement. $D(N)$ = count of distinct arrangements. Given $D(2)=3$, $D(10)=44499$, $D(20)=9204559704$. Find last 9 digits of $D(10000)$.
\section{Analysis}
\subsection*{Plane Partition Connection}

Each arrangement corresponds to a **plane partition** (3D Young diagram). The positions occupied by amoebas form an order ideal of the poset $\mathbb{N}^3$. The amoeba at $(x,y,z)$ requires that $(x-1,y,z)$, $(x,y-1,z)$, and $(x,y,z-1)$ are all occupied (or at the boundary).

\subsection*{Counting via Plane Partitions}

$D(N)$ equals the number of plane partitions of $N$ into parts arranged in a 3D staircase. More precisely, it counts the number of order ideals of size $2N+1$ in $\mathbb{N}^3$ reachable by sequential single-element extensions from $\{(0,0,0)\}$.

By a the
\section{Answer}

\[
\boxed{798443574}
\]

\end{document}
